\documentclass[lettersize,journal]{IEEEtran}

% ── Packages ──────────────────────────────────────────────────────────────
\usepackage[T1]{fontenc}
\usepackage[utf8]{inputenc}
\usepackage{amsmath,amsfonts}
\usepackage{amssymb}
\usepackage{bm}
\usepackage{algorithm}
\usepackage{algorithmic}
\usepackage{array}
\usepackage{textcomp}
\usepackage{stfloats}
\usepackage{url}
\usepackage{graphicx}
\usepackage{hyperref}
\usepackage{cite}
\usepackage{booktabs}
\usepackage{tabularx}
\usepackage[table]{xcolor}
\hyphenation{op-tical net-works semi-conduc-tor IEEE-Xplore}
% ── Placeholder colour ────────────────────────────────────────────────────
\newcommand{\placeholder}[1]{\textcolor{gray}{\textit{[#1]}}}

% ── Math shorthands ───────────────────────────────────────────────────────
% Preference vector and transformation matrix notation
\newcommand{\xvec}{\boldsymbol{x}}
\newcommand{\yvec}{\boldsymbol{y}}
\newcommand{\Ku}{K_{u}}
\newcommand{\Kuinv}{K_{u}^{-1}}

% ── Document ──────────────────────────────────────────────────────────────
\begin{document}

\title{Privacy-Preserving Intelligent Mobility for Individuals with
       Autism Spectrum Disorder}

\author{Ahmed Aouab Admou, Yahya Misbahi, Titouan Mathieu,
        Nawfal El-Hammadi, and Toumany Doumbouya%
\thanks{The authors are with the National Institute of Applied Sciences,
Toulouse, France.}
\thanks{Manuscript prepared for the Privacy and Intelligent Robotics
project, 2025--2026.}}

\markboth{Privacy-Preserving Intelligent Mobility for ASD}%
{Admou \MakeLowercase{\textit{et al.}}: Privacy-Preserving Intelligent Mobility for ASD}

\maketitle

% ─────────────────────────────────────────────────────────────────────────
\begin{abstract}
Smart mobility systems increasingly use artificial intelligence and
edge sensing to optimize urban transport, but they rarely account for
the sensory and behavioral needs of people with Autism Spectrum
Disorder (ASD). Existing sensory-aware navigation tools provide useful
personalization, yet they often rely on centralized storage or limited
privacy mechanisms for highly sensitive mobility and health-related
data. This study aimed to design and implement a privacy-preserving
intelligent mobility prototype for ASD users based on a three-tier
Cloud--MEC--Device architecture. The proposed system combines an
encrypted and matrix-obfuscated ASD profile vector, on-device
personalized inference, regional MEC context enrichment, and
personalized hierarchical federated learning. In the implemented
prototype, the phone-side model is decomposed as
$\mathbf{w}_i=\mathbf{w}_s+\mathbf{v}_i$, where only the shared
component $\mathbf{w}_s$ is federated and the personal component
$\mathbf{v}_i$ remains local. MEC nodes aggregate regional shared
updates and provide situational inputs such as crowding, delay, and
noise, while the cloud coordinates global shared model initialization.
Prototype validation using synthetic and manual demonstration data
shows that the web app performs local AI inference, shared updates are
federated through MEC nodes, personal components remain local, and the
cloud receives only regional shared model metadata. These results
demonstrate the technical feasibility of combining ASD-aware mobility
support with privacy-by-design learning.
\end{abstract}

\begin{IEEEkeywords}
Autism Spectrum Disorder, smart mobility, personalized federated
learning, mobile edge computing, privacy-by-design, sensory-aware
navigation.
\end{IEEEkeywords}

% ─────────────────────────────────────────────────────────────────────────
\section{Introduction}

\IEEEPARstart{S}{mart} city mobility systems increasingly rely on pervasive sensing,
context-aware services, and artificial intelligence to optimize urban
transport. These systems usually prioritize travel time, congestion, or
network efficiency, but they rarely model the sensory accessibility of
routes. For individuals with Autism Spectrum Disorder (ASD), mobility
is shaped not only by distance but also by noise, light, crowd density,
unexpected delays, and the availability of calm recovery spaces.
Exposure to these stimuli can produce anxiety, sensory overload, or
route avoidance, making conventional navigation insufficient for safe
and autonomous travel~\cite{maclennan2023}.

Several researchers have investigated sensory-aware and personalized
urban navigation. Mirri et al.~\cite{mirri2019} developed the PIUMA
application, which indexes places according to sensory characteristics
and enables route filtering from an individual user profile. Other
assistive mapping approaches have explored crowdmapping and inclusive
accessibility information. In parallel, privacy-preserving learning has
been studied for sensitive medical and mobility contexts. Federated
learning has been proposed to avoid centralized training data
collection, while differential privacy and edge computing have been
used to reduce the exposure of personal records~\cite{adnan2022,
alwarafy2020}. Recent reviews of Mobility-as-a-Service platforms have
also highlighted centralized cloud storage as a major privacy weakness
for mobility data~\cite{garroussi2025}.

However, previous work has not fully addressed the combination of
ASD-specific personalization, real-time regional mobility context, and
end-to-end privacy preservation. Existing systems often personalize
navigation by storing sensitive user profiles or by relying on
centralized service providers. They also rarely distinguish between
shared population-level learning and strictly personal adaptation. As a
result, it remains unclear how an ASD-aware mobility system can provide
personalized discomfort prediction while preventing raw sensory
profiles, physiological data, and personal model components from
leaving the user's device.

This study aimed to design and implement a privacy-preserving
intelligent mobility prototype for ASD users. The system represents
sensory sensitivities as a normalized profile vector, protects this
vector through user-specific matrix obfuscation and encryption, and
combines it locally with contextual and physiological features for
discomfort prediction. The learning architecture follows a
Cloud--MEC--Device hierarchy: the phone performs local inference and
personalized training, MEC nodes aggregate regional shared model
updates and provide situational context, and the cloud coordinates
global shared model initialization.

The main contribution is a fully functional prototype that integrates
personalized federated learning into this architecture. The phone model
is decomposed into a shared component and a private personal component,
so that $\mathbf{w}_i=\mathbf{w}_s+\mathbf{v}_i$. Only
$\mathbf{w}_s$ is sent to the MEC; $\mathbf{v}_i$ remains local. The
prototype also demonstrates MEC-based context enrichment, encrypted
profile synchronization, regional and global shared model endpoints,
and an operational web interface that runs local AI inference. The rest
of the paper reviews related work, presents the mathematical and
privacy framework, describes the implementation stack, and reports
prototype validation results obtained from synthetic and manual test
data.

% ─────────────────────────────────────────────────────────────────────────
% ─────────────────────────────────────────────────────────────────────────
\section{Related Work}
\label{sec:background}

\subsection{Inclusive Mobility and the Privacy Paradox}
Smart city mobility research increasingly uses sensing, artificial
intelligence, and real-time transport data to optimize urban flow.
However, inclusive mobility systems often still assume neurotypical
sensory processing. For many people with Autism Spectrum Disorder
(ASD), public space is experienced as a dense sensory environment in
which noise, light, crowding, smell, heat, route disruption, and lack of
predictable recovery spaces can directly affect autonomy and safety
~\cite{maclennan2023}. Personalized assistance therefore requires
fine-grained data such as sensory preferences, physiological stress
signals, behavioral indicators, and location context. This creates a
privacy paradox: the same information needed to adapt mobility support
can reveal health status, daily routines, and vulnerable places if it is
centralized or shared with third parties.
The relevant sensory barriers are multidimensional: auditory overload
from sirens or announcements, visual discomfort from flickering or
reflective light, social stress from dense crowds, olfactory triggers,
spatial disorientation in complex stations, and thermal discomfort in
poorly ventilated transit environments. Prior work also emphasizes the
importance of predictable routes and recovery spaces, since unexpected
events such as a cancelled bus or a closed passage can be more
disruptive than a longer travel time~\cite{maclennan2023,beck2025}.

\subsection{Sensory-Aware Navigation Systems}
Existing sensory-aware systems show that neurodivergent mobility cannot
be reduced to shortest-path routing. PIUMA indexes places according to
sensory features and lets users filter routes using an individual
profile~\cite{mirri2019}. Related crowdmapping approaches such as
mPASS, Maps4all, and CARES extend accessibility information by
collecting user reports about architectural or sensory barriers.
Together, these systems demonstrate the value of explicit sensory
information, visual route interfaces, and calm-space identification.
Their limitations are also consistent across the literature: coverage
depends heavily on community contributions, real-time events such as
delays or crowding are difficult to capture, and municipal accessibility
frameworks still prioritize physical disability over cognitive and
sensory accessibility~\cite{garroussi2025}.
PIUMA is particularly relevant because it combines three ideas that also
shape our prototype: a user-provided sensory profile, crowdmapped
ratings of sensory points of interest, and assistive routing based on
the user's current internal state. Earlier systems such as mPASS focus
more strongly on physical barriers, while Maps4all and CARES show that
community reporting and social participation are important but remain
insufficient when real-time regional context is missing.

\subsection{Privacy-Preserving AI for Mobility}
Most Mobility-as-a-Service platforms rely on centralized cloud
architectures, which expose sensitive records during transmission,
storage, analytics, or partner sharing~\cite{garroussi2025}. Standard
consent mechanisms may also be inaccessible to neurodivergent users
when they rely on legal language, long forms, or interface patterns that
reduce genuine understanding~\cite{beck2025}. Privacy-by-design
approaches therefore move computation closer to the user and minimize
what leaves the device. Multi-access Edge Computing (MEC) reduces
latency and limits raw-data exposure by processing regional context
near users~\cite{alwarafy2020}. Federated Learning (FL) trains models
locally and exchanges model updates rather than raw records. Differential
privacy (DP), usually implemented through clipping and calibrated noise,
can additionally harden FL updates against inference attacks, but it is
not the primary privacy mechanism in our prototype. These methods
support a decentralized architecture, but conventional FL still tends to
merge personal and shared learning into a single model.
Additional privacy mechanisms are needed around the user profile
itself. In our setting, the ASD sensitivity vector is not only encrypted
for storage but also transformed by a user-specific matrix derived from
credential-based key material. This matrix transformation is not treated
as a replacement for encryption; rather, it is an extra obfuscation step
that makes a stored vector harder to interpret if database records are
exposed. Edge processing complements this protection by keeping
physiological signals and stress inference close to the phone, where
immediate route adaptation can happen without broadcasting the user's
raw state to the cloud.

\begin{table}[!t]
\caption{Related Work Synthesis and Remaining Gap\label{tab:related_work_synthesis}}
\centering
\scriptsize
\renewcommand{\arraystretch}{1.2}
\begin{tabularx}{\columnwidth}{
>{\raggedright\bfseries}p{1.7cm}
>{\raggedright}X
>{\raggedright\arraybackslash}X}
\toprule
\rowcolor{gray!15}
\textbf{Area} & \textbf{Main Contribution} & \textbf{Remaining Gap} \\
\midrule
Sensory maps &
Expose noise, crowding, and other sensory attributes &
Often rely on static or crowdmapped data with limited real-time context \\
\midrule
MaaS privacy &
Identifies centralized storage and sharing risks &
Rarely addresses ASD-specific profile sensitivity and consent needs \\
\midrule
MEC and FL &
Moves computation closer to users and avoids raw-data upload &
Usually aggregates a single model without separating personal adaptation \\
\midrule
Accessible consent &
Improves readability and user autonomy &
Must be integrated into the technical privacy architecture \\
\bottomrule
\end{tabularx}
\end{table}

\subsection{Positioning of This Work}
The reviewed literature motivates a system that combines sensory-aware
mobility support, accessible privacy practices, and decentralized AI.
The gap addressed in this work is the lack of an integrated prototype
that separates regional knowledge from private personal adaptation. Our
architecture combines encrypted profile synchronization, additional
matrix obfuscation of the ASD profile vector, MEC-based situational
context, and personalized hierarchical FL. Unlike centralized mobility
platforms, raw sensory profiles and physiological context remain on the
phone. Unlike standard FL, the phone model is decomposed into a shared
component federated through MEC nodes and a private component that never
leaves the device.
The MEC layer also addresses a limitation of crowdmapped sensory
navigation by injecting situational data, such as metro delay, platform
crowding, construction noise, or quiet-area availability, into the local
phone-side model. Thus, the system learns from regional infrastructure
conditions without requiring the MEC to see the user's private ASD
profile or personal adaptation vector.
% ---------------------------------------------------------------
\section{System Architecture and Mathematical Framework}
% ---------------------------------------------------------------
\label{sec:architecture}
\subsection{Overview of the Three-Tier Hierarchical Architecture}

The proposed system is organized into three computational tiers that
interact bidirectionally: the user device (smartphone), Multi-access
Edge Computing (MEC) nodes deployed at transport infrastructure
locations, and a central cloud server. This hierarchy is denoted
Cloud~$\leftrightarrow$~MEC~$\leftrightarrow$~Device. Each tier is
responsible for a distinct scope of learning and inference. The device
tier performs real-time, personalized discomfort prediction and local
model fine-tuning using data that never leaves the hardware. The MEC
tier aggregates model weights from multiple devices within a geographic
region and specializes the shared model to local environmental patterns
such as infrastructure layout, typical crowd dynamics, and regional
climate. The cloud tier aggregates updates from all MEC nodes to
maintain a global model capturing broad ASD mobility behaviors across
the entire user population, and distributes generalized model
initializations back down to MEC nodes and devices. No raw user data
propagates beyond the device at any point in this pipeline.

The hierarchy also separates time scales. The phone reacts at the
shortest time scale, because discomfort prediction must occur during a
walk, transit change, or station arrival. MEC nodes operate at an
intermediate time scale, updating regional shared models and situational
context as local transport conditions evolve. The cloud operates at the
longest time scale, learning general initialization weights across
regions and users. This separation is important for ASD mobility because
the system must combine stable individual preferences with fast changes
in the environment, such as a delayed metro line, a crowded platform, or
unexpected construction noise.

Operationally, the device owns the complete personal context. It
combines the recovered ASD survey profile, current check-in answers,
optional wearable signals, and MEC-provided context into the local
feature vector. The MEC owns regional context and regional aggregation
only; it does not need to know which sensory dimensions are most
sensitive for a specific user. The cloud owns global coordination and
encrypted synchronization metadata. This division makes the privacy
boundary explicit and avoids treating the cloud as a trusted holder of
raw disability-related records.

\subsection{Threat Model and Privacy Boundary}

The system assumes an honest-but-curious cloud backend and honest-but-
curious MEC operators: these entities execute the protocol but may
attempt to infer sensitive information from the data they receive. The
network is assumed to be observable by an external adversary, so
transport encryption is required for deployment. The database may also
be compromised, motivating encryption of stored profile material and
obfuscation of profile vectors prior to persistence. Malicious clients
that deliberately submit poisoned model updates are considered a future
robustness problem and are outside the current prototype scope.

\begin{table}[!t]
    \centering
    \caption{Privacy Boundary of the Cloud--MEC--Device Architecture\label{tab:privacy_boundary}}
    \scriptsize
    \renewcommand{\arraystretch}{1.25}
    \begin{tabularx}{\columnwidth}{
        >{\raggedright\bfseries}p{1.5cm}
        >{\raggedright}X
        >{\raggedright\arraybackslash}X}
        \toprule
        \rowcolor{gray!15}
        \textbf{Layer} & \textbf{Visible Data} & \textbf{Hidden Data} \\
        \midrule
        Phone &
        ASD profile, raw check-in context, physiological signals,
        $\mathbf{v}_i$, and effective model $\mathbf{w}_i$ &
        Other users' profiles and model states \\
        \midrule
        MEC &
        Situational context, pseudonymous risk scores, shared
        updates $\mathbf{w}_s$, sample counts, regional aggregate model &
        ASD profile vector, physiological raw data, local dataset,
        personal component $\mathbf{v}_i$ \\
        \midrule
        Cloud &
        Encrypted profile records, derived check-in metrics, regional
        shared models, global shared model metadata &
        Raw profile vector, raw check-in payload, raw physiology,
        personal component $\mathbf{v}_i$ \\
        \bottomrule
    \end{tabularx}
\end{table}

\subsection{ASD Profile Representation and Input Feature Engineering}

Each user's sensory profile is represented as a normalized vector
\begin{equation}
    \mathbf{x} = (x_1, x_2, \dots, x_n)^\top \in [0,1]^n
\end{equation}
where each component $x_i$ encodes a scalar sensitivity score for a
distinct sensory or behavioral dimension, including noise tolerance,
light sensitivity, crowd density tolerance, susceptibility to routine
disruption, and transition difficulty. This vector is derived from a
structured onboarding questionnaire administered at first launch. The
vector $\mathbf{x}$ is computed locally and persists only on the device
in encrypted form, as described in Section~\ref{sec:crypto}.

At inference time, the profile vector is combined with a real-time
contextual feature vector
\begin{equation}
    \mathbf{c} = (c_1, c_2, \dots, c_m)^\top \in \mathbb{R}^m
\end{equation}
where contextual features include ambient noise level, estimated crowd
density at the current location, weather conditions, heart rate acquired
from a paired wearable device, time of day, and current transit mode.
The composite input to the on-device model is the concatenated feature
vector
\begin{equation}
    \mathbf{z} = [\mathbf{x}^\top \mid \mathbf{c}^\top]^\top
    \in \mathbb{R}^{n+m}
\end{equation}
The on-device model $f_{\boldsymbol{\theta}}$ maps $\mathbf{z}$ to a
predicted discomfort score
\begin{equation}
    d = f_{\boldsymbol{\theta}}(\mathbf{z}) \in [0, 1]
\end{equation}
where $d$ represents the estimated probability of an imminent sensory
overload or stress episode. The model architecture is a compact
feedforward neural network composed of a fully connected layer, a ReLU
activation, and a final classification layer, selected to remain
suitable for local simulation and on-device execution.

In the implemented prototype, the feature groups are explicitly
separated by origin. User-derived features include the ASD sensory
profile, current energy, breathing state, communication difficulty, and
need for isolation. Physiological features include heart rate, heart
rate variability when available, sleep quality, and wearable
availability. Local environmental features include perceived noise,
light, crowd density, routine disruption, transition difficulty, social
load, and calm-space availability. MEC situational features include
regional signals such as metro delay, transport incident state,
platform crowding, construction noise, and quiet-area availability.
Observable support signals include pacing, increased stimming, and
shutdown indicators. All these features are fused locally on the phone;
the MEC supplies only regional context and never receives the complete
personal feature vector.

The feature design is intentionally modular. Survey-derived features
change slowly and describe stable sensitivities; physiological and
check-in features change at session time; MEC situational features
change with the transport environment. Keeping these groups explicit
allows the prototype to support partial sensing. For instance, if a
wearable heart-rate signal is unavailable, the local model can still
combine the ASD profile, self-reported state, and MEC context. This
design also supports future ablation studies, because each feature group
can be removed or perturbed to measure its contribution to discomfort
prediction.

\subsection{Personalized On-Device Federated Learning}

To avoid forcing every user in a region to share a single identical
model, the phone model is decomposed into a shared component and a
personal adaptation component. For user $i$, the effective inference
model is
\begin{equation}
    \mathbf{w}_i = \mathbf{w}_s + \mathbf{v}_i
\end{equation}
where $\mathbf{w}_s$ is the shared regional component received from the
nearest MEC node and $\mathbf{v}_i$ is the private personal component
stored only on the user's phone. For a new user, $\mathbf{v}_i$ is
initialized to the zero vector, so the first effective model is simply
the MEC initialization.

During local training, the loss is computed on the effective model
$\mathbf{w}_s+\mathbf{v}_i$, but the two components are updated with
separate learning rates:
\begin{equation}
    \mathbf{w}_s \leftarrow \mathbf{w}_s
    - \eta_s \nabla \mathcal{L}_i(\mathbf{w}_s + \mathbf{v}_i)
\end{equation}
\begin{equation}
    \mathbf{v}_i \leftarrow \mathbf{v}_i
    - \eta_p \nabla \mathcal{L}_i(\mathbf{w}_s + \mathbf{v}_i)
\end{equation}
where $\eta_s$ controls how quickly the shared component adapts and
$\eta_p$ controls user-specific personalization. This design lets the
phone benefit from regional learning while preserving a private
adaptation vector that captures idiosyncratic sensory reactions,
feedback history, and behavioral patterns.

Only the updated shared component $\mathbf{w}_s$ is transmitted to the
MEC node. The personal component $\mathbf{v}_i$ is never serialized,
published, or stored outside the phone. The main privacy boundary
therefore comes from local data processing, non-transmission of raw
records, encrypted profile storage, and the strict separation of
$\mathbf{v}_i$. As an optional hardening layer, the shared update can be
clipped and noised before transmission:
\begin{equation}
    \tilde{\mathbf{w}}_{s,i} =
    \mathrm{clip}\!\left(\mathbf{w}_{s,i},\,C\right)
    + \mathcal{N}\!\left(\mathbf{0},\,\sigma^2\mathbf{I}\right)
\end{equation}
The transmitted payload therefore contains only
$\tilde{\mathbf{w}}_{s,i}$, the local sample count $n_i$, and coarse
training metadata. The raw dataset $\mathcal{D}_i$ and the private
component $\mathbf{v}_i$ never leave the device.

\subsection{Hierarchical Federated Aggregation}

\textbf{MEC-level aggregation.}
A MEC node $j$ collects only perturbed shared components from all
$K_j$ devices within its coverage area. It computes a weighted average
of the received shared models proportional to each device's local
sample size:
\begin{equation}
    \mathbf{w}_{s,\mathrm{region},j} =
    \frac{\displaystyle\sum_{i=1}^{K_j} n_i\,\tilde{\mathbf{w}}_{s,i}}
         {\displaystyle\sum_{i=1}^{K_j} n_i}
\end{equation}
The MEC node rejects any payload containing a personal component and
uses $\mathbf{w}_{s,\mathrm{region},j}$ as the updated regional
initialization distributed back to nearby phones.

\textbf{Cloud-level aggregation.}
The central cloud server receives only regional shared models from all
$J$ MEC nodes and performs a second-level weighted aggregation
proportional to the total number of device-samples contributed per MEC
node:
\begin{equation}
    \mathbf{w}_{s,\mathrm{cloud}} =
    \frac{\displaystyle\sum_{j=1}^{J} N_j\,\mathbf{w}_{s,\mathrm{region},j}}
         {\displaystyle\sum_{j=1}^{J} N_j}
\end{equation}
where $N_j = \sum_{i \in j} n_i$ is the total number of samples
aggregated under MEC node $j$. The updated global shared component is
broadcast to MEC nodes, which use it to initialize or refresh regional
models.

\textbf{Region switching.}
When a user moves from one MEC region to another, only the shared
component is replaced. For example, if a phone moves from Toulouse to
Paris, the effective model changes from
$\mathbf{w}_{s,\mathrm{Toulouse}}+\mathbf{v}_i$ to
\begin{equation}
    \mathbf{w}_i =
    \mathbf{w}_{s,\mathrm{Paris}} + \mathbf{v}_i
\end{equation}
The private vector $\mathbf{v}_i$ is preserved exactly. This prevents
loss of personalization while allowing the user to benefit from the new
regional MEC model.

\subsection{Personalized Hierarchical FL Protocol}

The complete learning procedure implemented by the prototype is
summarized as Algorithm~\ref{alg:pfl}. The algorithm highlights the
initialization chain, the local personalization step, and the fact that
only shared components are aggregated upstream.

\begin{algorithm}[!t]
\caption{Personalized hierarchical federated learning}
\label{alg:pfl}
\begin{algorithmic}[1]
\STATE \textbf{Input:} global shared model
$\mathbf{w}_{s,\mathrm{cloud}}$, MEC region $r$, phone user $i$, local
dataset $\mathcal{D}_i$
\STATE Cloud maintains $\mathbf{w}_{s,\mathrm{cloud}}$ from regional
shared submissions.
\STATE When MEC $r$ starts, request $\mathbf{w}_{s,r}$ from the cloud.
\IF{no regional model exists}
    \STATE $\mathbf{w}_{s,r} \leftarrow
    \mathbf{w}_{s,\mathrm{cloud}}$
\ENDIF
\STATE Phone requests $\mathbf{w}_{s,r}$ from the nearest MEC, then from
the cloud fallback if the MEC is unavailable.
\STATE For a new phone user, set $\mathbf{v}_i \leftarrow \mathbf{0}$.
\STATE Compute $\mathbf{w}_i=\mathbf{w}_{s,r}+\mathbf{v}_i$ for local
inference using profile, physiology, check-in, and MEC context.
\STATE Update $\mathbf{w}_{s,r}$ and $\mathbf{v}_i$ locally with
separate learning rates on
$\nabla\mathcal{L}_i(\mathbf{w}_{s,r}+\mathbf{v}_i)$.
\STATE Transmit only $\tilde{\mathbf{w}}_{s,i}$ and $n_i$ to the MEC.
\STATE MEC rejects any payload containing $\mathbf{v}_i$ or raw data.
\STATE MEC computes $\mathbf{w}_{s,\mathrm{region},r} =
\sum_i(n_i/N)\tilde{\mathbf{w}}_{s,i}$.
\STATE MEC uploads only the regional shared model and aggregate metadata.
\STATE On region switch, replace $\mathbf{w}_{s,r}$ with
$\mathbf{w}_{s,r'}$ and preserve $\mathbf{v}_i$.
\end{algorithmic}
\end{algorithm}

\subsection{Profile Vector Obfuscation}

To allow cloud-side synchronization of the ASD profile across user
devices without exposing the raw sensitivity vector, a user-specific
linear obfuscation is applied before any data leaves the device. A
transformation matrix $\mathbf{M}_u \in \mathbb{R}^{n \times n}$ is
constructed deterministically from key material derived from the user's
authentication credentials via a Key Derivation Function. The stored
representation of the profile is
\begin{equation}
    \mathbf{y} = \mathbf{M}_u \cdot \mathbf{x}
\end{equation}
The matrix $\mathbf{M}_u$ is never persisted; it is recomputed on
demand from the credentials at authentication time. Profile
reconstruction on a new device proceeds as
\begin{equation}
    \mathbf{x} = \mathbf{M}_u^{-1} \cdot \mathbf{y}
\end{equation}
For this reconstruction to be well-defined, $\mathbf{M}_u$ must be
invertible. This is enforced by construction: the KDF output seeds a
deterministic pseudorandom generator that produces a full-rank matrix.
The obfuscated vector $\mathbf{y}$ is subsequently encrypted before
storage as described in Section~\ref{sec:crypto}. This matrix
transformation is therefore not intended to replace encryption, but to
act as an additional obfuscation layer. If an attacker observes a stored
numeric vector or an intermediate serialized representation, the values
do not directly correspond to interpretable sensitivities such as
noise, light, crowd density, or routine disruption. Without the
credential-derived matrix, the transformed coordinates are difficult to
interpret semantically even before considering the cryptographic
protection applied at rest.

\subsection{Cryptographic Data Protection}
\label{sec:crypto}

The system employs a two-layer key hierarchy to protect stored data. A
Data Encryption Key (DEK) is generated per user as a cryptographically
random symmetric key. The obfuscated profile vector $\mathbf{y}$ is
encrypted under AES-256-GCM using the DEK:
\begin{equation}
    C_{\mathbf{y}} =
    \mathrm{Enc}_{\mathrm{AES\text{-}256\text{-}GCM}}(\mathrm{DEK},\;\mathbf{y})
\end{equation}
The DEK is itself encrypted under a Key Encryption Key (KEK) derived
from the user's password and a random salt via Argon2id:
\begin{equation}
    \mathrm{KEK} = \mathrm{KDF}_{\mathrm{Argon2id}}(\mathrm{password},\;\mathrm{salt})
\end{equation}
\begin{equation}
    C_{\mathrm{DEK}} =
    \mathrm{Enc}_{\mathrm{AES\text{-}256\text{-}GCM}}(\mathrm{KEK},\;\mathrm{DEK})
\end{equation}
Only $C_{\mathbf{y}}$ and $C_{\mathrm{DEK}}$ are stored in the
database alongside the salt. The password and plaintext DEK are never
persisted. At authentication, the KEK is recomputed from the submitted
password and the stored salt, the DEK is decrypted transiently in
memory, and the profile is decrypted for use. On devices equipped with
hardware-backed secure storage such as Apple Secure Enclave, Android
StrongBox, or a Trusted Platform Module, the KEK or DEK may
alternatively be bound to hardware-attested key material, providing an
additional layer of protection against software-level extraction.

This mechanism supports cross-device recovery of the questionnaire
profile but not automatic recovery of all phone-local learning state.
When a user authenticates on a new device, the encrypted database record
can be used to recover the obfuscated and encrypted ASD survey profile,
provided that the user supplies the correct credentials. The personal
model component $\mathbf{v}_i$, however, is designed as a phone-local
adaptation vector and is not synchronized in the current prototype. This
choice protects personalization from server-side exposure but creates a
portability limitation: a new device can reconstruct the user's survey
profile and initialize $\mathbf{w}_s$ from the MEC or cloud, but it
starts with a fresh local personal vector unless a future secure
transfer mechanism is added.

One possible extension is to apply the same protection pattern to
$\mathbf{v}_i$ when cross-device transfer is explicitly requested by the
user. In such a design, the personal vector could be transformed with a
credential-derived matrix, encrypted with a DEK, and stored as an
optional encrypted recovery object. This would preserve the principle
that the server cannot interpret the personalized component while giving
users a controlled way to migrate local adaptation between trusted
devices.


% ---------------------------------------------------------------
\section{Implementation}
% ---------------------------------------------------------------

\subsection{Web Application and Phone-Side AI Prototype}

The current application prototype is implemented as a NiceGUI web
application in Python. Although it is delivered as a web interface for
demonstration purposes, it follows the same logical boundary as the
phone client: the user's active check-in context, personalized model
state, and personal adaptation vector remain in the client session. The
interface includes authentication, onboarding, manual check-in,
activity history, settings, and researcher/admin dashboards.

The check-in screen collects internal state, environmental conditions,
optional physiological signals, mobility context, and observable
support signals. These values are converted into a compact feature
vector $\mathbf{z}$ and evaluated by a lightweight local AI model. This
prototype model is implemented as a linear personalized model with a
sigmoid output so that the Personalized Federated Learning mechanism can
be executed directly inside the web app without requiring a heavy
runtime. The effective model is $\mathbf{w}_i=\mathbf{w}_s+\mathbf{v}_i$:
the shared component $\mathbf{w}_s$ is initialized from the MEC
regional model, and the private component $\mathbf{v}_i$ is initialized
to zero for new users. After each check-in, the local training step
updates both components, but only the shared component is published to
the MEC broker. The personal vector remains in the web session and is
not included in MQTT or REST payloads.

The local feature vector may also include situational context obtained
from the currently associated MEC node. For example, when the MEC
reports that a metro line is delayed, that a station is crowded, or
that a platform is noisy, these values are treated as contextual input
features on the phone. The MEC does not need access to the user's ASD
profile or physiological signals to provide this context; it only
supplies regional environmental information that improves the local
prediction.

In the prototype, this behavior is visible through the web interface.
The check-in workflow produces a discomfort estimate immediately after
the user submits the local context. The model state stored in the client
session records the shared source, regional identifier, training round,
and private adaptation vector. These fields allow the interface to show
whether the shared component came from the MEC, from a cloud fallback,
or from a deterministic offline default. The same service also prepares
the shared update payload sent to the MEC, ensuring that $\mathbf{v}_i$
is excluded before serialization.

If the MEC API is reachable, the client obtains its initial shared
component from \texttt{/fl/model}. If the MEC is temporarily
unavailable, the client attempts to retrieve the regional shared model
from the cloud backend; if both services are unavailable, it falls back
to a deterministic default shared model so that inference remains
available offline. When the user later connects to a different MEC
region, the client replaces only
$\mathbf{w}_s$ with the new regional shared model and preserves
$\mathbf{v}_i$, maintaining personalization across region changes.

\subsection{Backend Server and Cloud Coordination Layer}

The central backend server is implemented in TypeScript with Express,
Prisma, and a SQLite database for the prototype deployment. It exposes a
REST API for authentication, encrypted profile synchronization,
privacy-preserving check-in storage, admin analytics, MEC node status,
and federated model coordination. The backend represents the cloud tier
of the hierarchy. It stores encrypted transformed ASD profile vectors,
encrypted DEKs, salts, check-in derived metrics, MEC metadata, regional
shared model states, and the global shared model state.

Authentication uses Express sessions. At login, the backend verifies the
password hash, reconstructs the credential-derived matrix, decrypts the
stored transformed profile, and recovers the normalized ASD vector
transiently for the authenticated session. The database stores only
encrypted profile material and encrypted DEKs; no plaintext ASD vector
is persisted. Check-ins are also reduced before persistence: the backend
stores derived metrics and prediction results rather than raw contextual
payloads.

For federated coordination, the backend exposes
\texttt{/api/federated/global-model},
\texttt{/api/federated/regional-model}, and
\texttt{/api/federated/updates}. A new MEC node requests a regional
model from the cloud. If no regional state exists, it is initialized
from the global shared model. MEC nodes later submit only regional
shared weights and sample counts to the cloud. The backend aggregates
these regional shared models into a global shared component and exposes
the resulting version for downstream initialization.

\subsection{MEC Nodes and Environment Simulation}

Each MEC node is deployed as a containerized microservices stack
orchestrated with Docker and Docker Compose. The stack comprises three
containers: \texttt{mec-mosquitto}, an Eclipse Mosquitto MQTT broker
responsible for message routing; \texttt{mec-api}, a FastAPI service
managing control logic, environmental context, and client
communication; and \texttt{mec-fl}, the federated aggregation service
implementing FedAvg in NumPy.

Real-time communication between devices and MEC nodes uses MQTT across
a structured topic hierarchy. The topic \texttt{tsa/prediction} carries
risk scores and zone identifiers. The topic \texttt{tsa/fl/update}
carries only shared model components $\mathbf{w}_s$ and sample sizes
submitted by devices for aggregation. The topic \texttt{tsa/context/*}
supports bidirectional environmental data exchange. The topic
\texttt{tsa/recommendation/\{id\}} delivers personalized routing advice
to individual clients.

The MEC layer also provides situational mobility context to the phone.
Examples include metro or bus delay, station crowding, construction
noise, platform congestion, incident reports, or the availability of a
quiet recovery area. These data describe the environment rather than
the individual user. They are sent from the MEC to the phone and are
then appended locally to the phone feature vector before inference. In
other words, the MEC enriches the local model input with regional
context, but the final discomfort prediction still runs on the phone
alongside the private ASD profile, physiological data, and personal
adaptation vector. This preserves the privacy boundary while allowing
the model to react to events such as a delayed metro, an unusually dense
station, or a noisy platform.

A key implementation constraint addressed in the \texttt{mec-fl}
service is enforcing the personalized FL privacy boundary. The service
accepts \texttt{shared\_weights} or \texttt{sharedWeights} fields but
rejects any payload containing a personal component. The received shared
model vectors are reconstructed from JSON lists and aggregated using the
weighted FedAvg computation defined in~(10). After aggregation, the MEC
updates its local \texttt{/fl/model} endpoint, publishes the new
regional shared model to subscribed phones, and forwards only aggregate
regional metadata to the cloud backend.

Since physical sensor infrastructure was not available for evaluation,
a Station Controller module was developed to simulate real-time
environmental conditions at transit stations. This simulator generates
synthetic streams of noise level, crowd density, and service delay
measurements published to the \texttt{tsa/context/*} MQTT topics and
consumed by MEC node services and mobile clients subscribed to their
current location.

\subsection{Federated Learning Prototype and Cloud Coordination}

A standalone federated learning prototype was implemented in Python
using PyTorch and Flower to validate the privacy-preserving learning
workflow independently of the deployed MEC stack. Three model roles
were implemented. The first is \texttt{PhoneModel}, trained locally on
each simulated device using private synthetic data representing
mobility-related and sensory-related preferences. The second is
\texttt{GlobalModel}, which shares the same architecture as
\texttt{PhoneModel} and is aggregated by the Flower server via FedAvg.
The third is \texttt{MECModel}, a lightweight edge model designed to
operate on non-identifying aggregate summaries rather than raw user
profiles.

On the client side, the Flower \texttt{NumPyClient} interface implements
local personalized training and submission of protected shared updates
to the server. Parameter clipping and Gaussian noise injection are
included as optional DP-style hardening for experiments where update
inference risk must be reduced.
The implementation mirrors the deployed web-app logic: each simulated
client keeps a private $\mathbf{v}_i$ initialized to zero, computes
gradients on $\mathbf{w}_s+\mathbf{v}_i$, updates both local
components, and returns only the clipped/noised shared component. Only
coarse training metrics are exposed to the server; no user identifier,
raw questionnaire value, or personal component is transmitted. The
server coordinates aggregation rounds using the \texttt{FedAvg}
strategy and updates the global shared model without accessing any
private local dataset. The \texttt{MECModel} establishes the code
structure required for future edge-side learning and local decision
support. Current training data is synthetic, defining the input feature
structure required for future integration with the onboarding
questionnaire and sensor streams.

The full system coordinates four implementation environments:
TypeScript/Express for the cloud backend, Python/NiceGUI for the
interactive client prototype, Python/FastAPI and MQTT services for MEC
nodes, and PyTorch/Flower for the federated learning simulation.
Inter-service communication uses MQTT for low-latency phone--MEC
updates and REST for account management, encrypted profile recovery,
cloud model initialization, and upstream regional model synchronization.

\subsection{Operational Stack and Startup Flow}

The project repository is organized into four main folders:
\texttt{backend}, \texttt{ui}, \texttt{mec}, and \texttt{ml}. The
\texttt{backend} folder contains the Express/Prisma cloud API. The
\texttt{ui} folder contains the NiceGUI web application that simulates
the phone client and runs the personalized model locally. The
\texttt{mec} folder contains the Docker Compose stack for edge
execution, including Mosquitto, the MEC API, the context service, and
the regional FL aggregator. The \texttt{ml} folder contains the
Flower/PyTorch learning prototype used to validate the training logic
independently from the live web interface.

A root launcher script, \texttt{run\_project.ps1}, starts the backend,
the UI, and the MEC Docker stack when Docker is available. The startup
sequence follows the learning hierarchy. First, the backend exposes the
global shared model. Second, a new MEC node starts and requests its
regional model from the backend. If the backend has no regional state
for that MEC, the node initializes from the global shared component.
Third, a phone session starts and requests the current shared model from
the nearest MEC. If the MEC is not reachable, the phone requests the
regional shared model from the cloud backend. If neither service is
reachable, it uses a local default shared component so that inference
continues in degraded mode. When the MEC becomes reachable again, the phone can refresh
$\mathbf{w}_s$ while keeping $\mathbf{v}_i$ unchanged.

\begin{table*}[!t]
    \centering
    \caption{Main Communication Channels and Privacy Role\label{tab:protocol}}
    \scriptsize
    \renewcommand{\arraystretch}{1.2}
    \begin{tabularx}{\textwidth}{
        >{\raggedright\bfseries\arraybackslash}p{3.6cm}
        >{\raggedright\arraybackslash}p{3.0cm}
        >{\raggedright}X
        >{\raggedright\arraybackslash}X}
        \toprule
        \rowcolor{gray!15}
        \textbf{Channel} & \textbf{Direction} &
        \textbf{Payload} & \textbf{Privacy Role} \\
        \midrule
        \texttt{/api/federated/global-model} &
        Cloud $\rightarrow$ MEC &
        Global shared weights and version &
        Initializes new MEC nodes without user data \\
        \midrule
        \texttt{/api/federated/regional-model} &
        Cloud $\rightarrow$ MEC/Phone &
        Regional shared model or cloud fallback &
        Provides fallback when MEC is unreachable \\
        \midrule
        \texttt{/fl/model} &
        MEC $\rightarrow$ Phone &
        Regional shared component $\mathbf{w}_s$ &
        Initializes phone shared model only \\
        \midrule
        \texttt{tsa/fl/update} &
        Phone $\rightarrow$ MEC &
        Shared weights and sample count &
        Excludes $\mathbf{v}_i$ and raw data \\
        \midrule
        \texttt{tsa/context/*} &
        MEC $\leftrightarrow$ Phone &
        Situational mobility context &
        Adds local context without exposing profile \\
        \midrule
        \texttt{/api/federated/updates} &
        MEC $\rightarrow$ Cloud &
        Regional shared model and aggregate metadata &
        Enables cloud aggregation without phone records \\
        \bottomrule
    \end{tabularx}
\end{table*}


% ---------------------------------------------------------------
\section{Evaluation and Results}

\subsection{Evaluation Criteria}

The prototype evaluation focuses on technical feasibility rather than
clinical efficacy. Four criteria are considered. First, the local
inference path must remain available even when the MEC is unreachable,
because the user may still need support during connectivity loss.
Second, the federated path must enforce the privacy boundary by
transmitting only shared model components and aggregate metadata. Third,
regional adaptation must be observable: a MEC node should initialize
from the cloud, accept nearby shared updates, and expose a regional
model to phones. Fourth, personalization must survive region switching,
meaning that a new regional shared component can be loaded without
overwriting the private vector $\mathbf{v}_i$.

These criteria correspond to the main claims of the architecture. They
also define the next empirical protocol. Once participant data are
available, the system should be compared against a cloud-only baseline,
a MEC-only shared model, a purely local model, and the proposed
personalized hierarchical model. Metrics should include prediction
accuracy, calibration of discomfort scores, latency, communication
volume, battery impact, and subjective usability. Privacy evaluation
should report the type of data visible at each tier, the amount of noise
added to shared updates, and the resulting privacy budget when a formal
accountant is integrated.

\subsection{Federated Learning Prototype Evaluation}

The federated learning prototype was evaluated locally using Flower's
simulation runtime. The objective of this evaluation was not to measure
final predictive performance on a real ASD mobility dataset, which was
not yet available, but to validate the execution of the complete
privacy-preserving training pipeline. The experiment verified that
simulated phone clients can train locally, protect their model updates,
transmit only perturbed parameters, and contribute to a global FedAvg
aggregation round.

Due to local memory constraints in the WSL execution environment, the
simulation was configured with two virtual supernodes. At each round,
one client was sampled for local training and evaluation. The server was
configured for three federated rounds, one local epoch per client, a
clipping norm of $C=1.0$, and Gaussian update noise with standard
deviation $\sigma=0.01$. The round timeout was increased to 120 seconds
to account for the overhead of the Flower and Ray simulation backend.

\begin{table}[!t]
    \centering
    \caption{Federated Learning Simulation Results Over Three Rounds\label{tab:fl_results}}
    \scriptsize
    \renewcommand{\arraystretch}{1.2}
    \begin{tabular}{ccc}
        \toprule
        \textbf{Round} & \textbf{Distributed Loss} & \textbf{Client Failures} \\
        \midrule
        1 & 0.7357 & 0 \\
        2 & 0.6883 & 0 \\
        3 & 0.6765 & 0 \\
        \bottomrule
    \end{tabular}
\end{table}

As shown in Table~\ref{tab:fl_results}, the distributed loss decreases
across the three rounds, from 0.7357 to 0.6765. This confirms that the
global model receives useful information from local client updates and
that the FedAvg aggregation loop executes correctly. No client failures
were observed in the final stabilized configuration. The experiment
therefore validates the functional feasibility of the federated
learning workflow in the proposed architecture.

The evaluation also confirmed the privacy-preserving data flow of the
prototype. Raw synthetic user records remained local to each simulated
phone client. The server received only clipped and noised model
parameters, together with coarse training metrics. This behavior is
consistent with the system objective of avoiding centralised exposure
of sensitive ASD-related mobility profiles.

The loss values should be interpreted as an engineering validation
signal rather than as a clinical performance claim. The goal of this
experiment was to verify that the training loop, parameter exchange,
aggregation, and privacy filters operate coherently. A future empirical
study will require participant-approved labels, repeated mobility
episodes, and comparison against baselines such as a non-personalized
global model, a regional-only model, and a purely local model. These
baselines will make it possible to quantify the contribution of
personalization, regional MEC context, and hierarchical aggregation.

\subsection{End-to-End Prototype Validation}

The implemented prototype was also validated as an integrated
Cloud--MEC--Device stack. The objective was to verify that each system
component executes its expected role and that the privacy boundaries in
Table~\ref{tab:privacy_boundary} are respected during normal operation.
The web application executes the local personalized model during
check-in, the MEC layer exposes regional shared-model initialization,
and the backend stores cloud and regional shared-model metadata.

\begin{table}[!t]
    \centering
    \caption{Functional Validation of the Implemented Prototype\label{tab:functional_validation}}
    \scriptsize
    \renewcommand{\arraystretch}{1.2}
    \begin{tabularx}{\columnwidth}{
        >{\raggedright\bfseries}p{2.0cm}
        >{\raggedright}X
        >{\raggedright\arraybackslash}p{1.7cm}}
        \toprule
        \rowcolor{gray!15}
        \textbf{Function} & \textbf{Validated Behavior} & \textbf{Result} \\
        \midrule
        Phone-local AI &
        Check-in features are evaluated by the personalized
        $\mathbf{w}_s+\mathbf{v}_i$ model in the web application &
        Functional \\
        \midrule
        Personalization privacy &
        The phone stores $\mathbf{v}_i$ locally and publishes only
        shared weights to \texttt{tsa/fl/update} &
        Functional \\
        \midrule
        MEC aggregation &
        The MEC aggregator accepts shared weights and rejects personal
        component fields &
        Functional \\
        \midrule
        Cloud coordination &
        The backend exposes global/regional initialization endpoints and
        accepts regional shared updates &
        Functional \\
        \midrule
        Region switching &
        The phone replaces only $\mathbf{w}_s$ when a new region model
        is received and preserves $\mathbf{v}_i$ &
        Functional \\
        \midrule
        Context enrichment &
        MEC situational context can be injected into local inference
        without transmitting the user's ASD profile &
        Functional \\
        \bottomrule
    \end{tabularx}
\end{table}

\subsection{Prototype Data and Demonstration Results}

Because no real participant study has yet been conducted, the current
data are prototype data generated by scripted phone clients, manual web
check-ins, and MEC environment simulations. These data exercise the
same communication and storage paths intended for deployment: check-ins
produce local discomfort predictions, phone clients generate shared
model updates, MEC nodes aggregate regional shared models, and the
cloud backend records global/regional model state. This provides
evidence of technical feasibility, but it should not be interpreted as
clinical validation of ASD stress prediction.

\begin{table}[!t]
    \centering
    \caption{Prototype Demonstration Results Using Synthetic and Manual Test Data\label{tab:demo_results}}
    \scriptsize
    \renewcommand{\arraystretch}{1.2}
    \begin{tabularx}{\columnwidth}{
        >{\raggedright\bfseries}p{2.2cm}
        >{\raggedright}X
        >{\raggedright\arraybackslash}p{1.5cm}}
        \toprule
        \rowcolor{gray!15}
        \textbf{Scenario} & \textbf{Observation} & \textbf{Outcome} \\
        \midrule
        Normal check-in &
        The web app generates a local stress estimate and logs derived
        metrics rather than raw context &
        Passed \\
        \midrule
        Delayed metro context &
        MEC situational delay is represented as local contextual input,
        increasing the model's available environmental evidence &
        Passed \\
        \midrule
        Shared update upload &
        Phone clients publish shared weights and sample counts; no
        personal component is included &
        Passed \\
        \midrule
        Regional aggregation &
        MEC computes weighted FedAvg over shared updates and republishes
        a regional shared model &
        Passed \\
        \midrule
        Cloud synchronization &
        Backend accepts regional shared model metadata and updates the
        global shared model version &
        Passed \\
        \bottomrule
    \end{tabularx}
\end{table}

% ---------------------------------------------------------------
\section{Discussion}
\subsection{Limitations of the Federated Learning Prototype}

The current federated learning evaluation remains preliminary. First,
the dataset used in the prototype is synthetic and only approximates the
structure of future ASD mobility features. It is sufficient to validate
the software pipeline, but not to draw conclusions about real-world
prediction accuracy. Second, differential privacy is treated as an
optional update-hardening mechanism rather than a core requirement of
the architecture; no formal privacy budget $(\varepsilon, \delta)$ has
yet been calibrated. Third, the MEC model is currently evaluated
separately from the Flower aggregation loop. Future work should
integrate MEC-level aggregation directly into the hierarchical
Cloud--MEC--Device workflow described in Section~\ref{sec:architecture}.

Despite these limitations, the prototype demonstrates that the proposed
architecture can support local training, protected update transmission,
and server-side aggregation without collecting raw user data.

A second limitation concerns device portability. The encrypted database
record allows an authenticated user to recover the ASD survey profile on
a new device, because the profile is stored as an encrypted and
matrix-obfuscated vector. The personalized model component
$\mathbf{v}_i$ is different: by design it remains local to the original
phone and is not uploaded to the MEC or cloud. Consequently, a new
device can initialize from the current MEC or cloud shared component and
recover the survey profile, but it cannot automatically recover the
previous device's learned personal adaptation. A future version could
make this optional by applying the same obfuscation and encryption
mechanisms to $\mathbf{v}_i$ before storing or transferring it, while
still ensuring that the server cannot interpret the personal vector.

\subsection{Security and Privacy Limitations}

The matrix transformation should be interpreted as an additional
obfuscation mechanism rather than a standalone cryptographic guarantee.
Encryption remains the primary confidentiality mechanism for stored
profile data. Similarly, differential privacy should be understood as an
optional protection for federated shared updates, not as the foundation
of privacy in the current design. The current prototype also does not
yet implement secure aggregation, Byzantine-robust aggregation, formal
client attestation, or a calibrated differential privacy accountant.
MQTT is used for local prototype communication and should be deployed
with TLS, authentication, and topic-level access control in a production
setting. Finally, the prototype data are synthetic or manually
generated; real-world validation requires an ethically approved study
with ASD participants or caregivers, carefully designed consent
materials, and controlled evaluation of prediction accuracy, latency,
usability, and safety.

% ---------------------------------------------------------------
\section{Conclusion}
% ---------------------------------------------------------------

This paper presented a privacy-preserving intelligent mobility
prototype for individuals with ASD. The system combines local
discomfort prediction, encrypted and matrix-obfuscated profile
synchronization, MEC-based situational context, and hierarchical
personalized federated learning. Its main design principle is that raw
sensory profiles, physiological signals, mobility context, and the
personal adaptation component of the model remain on the phone, while
only shared model updates are aggregated by MEC nodes and then by the
cloud.

The implemented prototype demonstrates the feasibility of this
Cloud--MEC--Device architecture. The web application performs local AI
inference using the effective model
$\mathbf{w}_i=\mathbf{w}_s+\mathbf{v}_i$, MEC nodes provide regional
context and aggregate shared updates, and the cloud maintains global and
regional shared model initialization. Region switching is also handled
without losing personalization, because only the shared component is
replaced while the private component is preserved. Prototype validation
with synthetic and manual demonstration data confirmed that the main
software paths execute as intended and that no raw personal data is
required for federated learning.

Future work will focus on training and evaluating the model with
ethically collected ASD mobility data, calibrating formal differential
privacy guarantees, strengthening secure aggregation and client
attestation, and measuring prediction accuracy, latency, usability, and
perceived safety in realistic mobility scenarios.


\section*{Acknowledgment}

The authors acknowledge the use of AI-based writing assistance only for
English formulation, grammar, and clarity improvements. The technical
content, system design, interpretation, and final manuscript decisions
were reviewed and validated by the authors, who remain fully responsible
for the submitted work.


% ---------------------------------------------------------------
\bibliographystyle{IEEEtran}
\begin{thebibliography}{1}

\bibitem{maclennan2023}
K.~MacLennan et al., ``Sensory Experiences in Public Spaces: A
Qualitative Study,'' \textit{Autism}, vol.~27, no.~4, 2023.

\bibitem{mirri2019}
S.~Mirri, C.~Prandi, and P.~Salomoni, ``How do sensory features of
places impact on spatial exploration of people with autism? A user
study,'' \textit{Information Technology and Tourism}, vol.~20,
pp.~1--18, 2019.

\bibitem{garroussi2025}
H.~Garroussi et al., ``A systematic review of data privacy in
Mobility-as-a-Service (MaaS),'' \textit{Transportation Research
Interdisciplinary Perspectives}, 2025.

\bibitem{beck2025}
S.~J.~Beck et al., ``Guidelines for the Creation of Accessible Consent
Materials and Procedures,'' \textit{Autism in Adulthood}, 2025.

\bibitem{adnan2022}
N.~Adnan et al., ``Federated Learning with Differential Privacy for
Medical Data,'' \textit{Scientific Reports}, vol.~12, no.~1, 2022.

\bibitem{alwarafy2020}
A.~Alwarafy, K.~A.~Al-Thelaya, M.~Abdallah, J.~Schneider, and
M.~Hamdi, ``A survey on security and privacy issues in edge computing
assisted Internet of Things,'' \textit{IEEE Internet of Things
Journal}, vol.~8, no.~6, pp.~4004--4022, 2020.

\end{thebibliography}

\end{document}
